\section{Interpretable Rationale Queries (L3)}

\subsection{Overview}
In this section and the next, we will explore queries that necessitate external data to furnish rationales for their resolution. These queries demand not only a grasp of the factual content but also the ability to comprehend and apply domain-specific \textbf{rationales} that are integral to the data's context. We classify these queries into two categories based on the nature of the rationales involved: queries based on interpretable rationales and those based on hidden rationales, as illustrated in the Figure ~\ref{fig:rationale}.

\MFUsentencecase{\thirdlevel} represent a relatively straightforward category within applications that rely on external data to provide rationales. The auxiliary data for these types of queries often include clear explanations of the thought processes used to solve problems. The data can be organized in several forms:

\begin{itemize}
    \item \textbf{Plain Texts}: Textual descriptions are the most common form of presenting interpretable rationales. These may include specialized or official documents such as handbooks or guidelines, as well as domain-specific manuals or operational guides. These texts articulate the reasoning processes that facilitate decision-making in complex scenarios. For example, documents such as FDA Guidance for pharmaceutical factories or medication guides for physicians provide insights into how experts, like FDA officers or doctors, approach specific cases.
    \item \textbf{Structured Instructions}: More explicit reasoning relationships or decision pathways might be presented in a structured format. These rationales can be understood as either a Text-Conditioned Moore Machine or a Text-Conditioned Mealy Machine. In the theory of computation, a Moore machine is a finite-state machine where the output values are determined solely by its current state\footnote{\url{https://en.wikipedia.org/wiki/Moore_machine}}. The conditions that control state transitions are often expressed in text, which LLMs need to interpret, unlike traditional programs that operate on native code. For instance, consider a customer supporting agent that follows a handbook to handle user's request to product changing or refunding. Similarly, a Mealy machine is a finite-state machine where output values are determined by both its current state and the inputs\footnote{\url{https://en.wikipedia.org/wiki/Mealy_machine}}. The distinction here is that actions (such as API calls) are determined not only by the state but also by the textual messages associated with transitions from the previous state. 
    Naturally, these domain-specific rationales can be represented in formats such as workflows, decision trees, or pseudocode.
\end{itemize}

Here are some examples of queries at this level:
\begin{itemize}
    \item \textit{How should a patient with chest pain and specific symptom descriptions be diagnosed and treated} (given a chest pain management guideline)
    \item \textit{How to respond to a user's question in a real-life scenario?} (given a customer service workflow)
\end{itemize}

\begin{figure}[t]
    \centering
    \includegraphics[width=0.8\textwidth]{contents/images/level34.png}
    \caption{Demonstration of Rationale Queries}
    \label{fig:rationale}
\end{figure}




\subsection{Challenges and Solutions}

In the realm of \thirdlevel, an additional challenge is integrating domain-specific rationales into LLMs in an comprehensible manner. The primary challenges are as follows:

\begin{itemize}
    \item \textbf{Prompt Optimization Costs}: The process of optimizing prompts is marked by high time and computational demands. Distinct queries demand tailored background knowledge and decision-making criteria, necessitating diverse examples. While manually designed prompts can be highly effective, they are labor-intensive and time-consuming. Furthermore, training models to generate tailored prompts for various queries incurs significant computational overhead.
    \item  \textbf{Limited interpretability}: The impact of prompts on LLMs is opaque. In many cases, access to the internal parameters of LLMs is typically restricted, complicating efforts to determine the impact of various prompts on these models. This lack of transparency hinders our ability to consistently understand and verify the interpretability of LLM responses to different prompts.
\end{itemize}

\subsection{Prompt Tuning}

For \thirdlevel, the key issue is how to effectively integrate rationales provided by external data into LLMs and ensure that these models can accurately follow and react based on these rationales. Text2MDT~\cite{zhu2024text2mdt} offers a viable demonstration, introducing two methods for automatically extracting medical decision trees from medical guidelines and textbooks. This process clarifies the logical chains within lengthy medical texts, making them more comprehensible. Similarly, MedDM~\cite{li2023meddm} has developed a format for clinical guidance trees that can be executed by LLMs, proposing a methodology for reasoning on these executable CGTs and a framework for multi-turn dialogues between patients and LLMs. InstructRec~\cite{zhang2023recommendation} aims to leverage the capabilities of LLMs in recommendation systems, designing a universal format to describe a user's preferences, intentions, task forms, and context using natural language, thereby creating a high-performing, language-based recommendation system.

Integrating rationales directly as natural language instructions into LLMs does not necessarily yield optimal performance, and manually designing prompts can be time-consuming. To address this, the employment of prompt tuning techniques becomes essential to enhance the LLMs' capability to adhere to specific rationales. One effective methodology is the application of reinforcement learning, as evidenced by the TEMPERA framework~\cite{zhang2022tempera}, which designs prompts incorporating limited instructions, examples, and verbalizers within the action space of reinforcement learning. Here, the LLM's probability of generating correct responses serves as the reward, guiding the model to discover the optimal prompt configuration across datasets. Similarly, Rlprompt~\cite{deng2022rlprompt} adopts a reinforcement learning approach, training an adaptor to assist smaller language models in producing optimal prompts based on feedback concerning the relative accuracy of LLM responses. Another innovative strategy, Directional Stimulus Prompting, leverages the performance of LLMs on downstream tasks as a reward mechanism. This method trains models to extract and utilize directional stimuli—specific cues or keywords tailored to individual instances—as prompts, thereby ensuring the LLMs' actions align more closely with expected outcomes.

Additionally, for optimization within discrete prompt spaces, edit-based methodologies such as GrIPS~\cite{prasad2022grips} are utilized. This technique involves using a small dataset as a scoring set to experiment with various prompt modifications—including deletions, swaps, paraphrases, and additions—to ascertain the most effective prompt configurations swiftly and effectively.

Recent advancements~\cite{zhou2022large, pryzant2023automatic} have also seen the rise of using LLMs themselves to facilitate prompt optimization. OPRO~\cite{yang2024large_1} employs an LLM both to generate new prompt solutions based on historical data and their associated performance metrics and to score these prompts, thus streamlining the optimization process. Furthermore, the Reflexion framework~\cite{shinn2024reflexion} introduces a novel prompt optimization approach based on linguistic feedback, using a language model to analyze and store reflections on LLM outputs in an episodic memory buffer. This memory component aids in refining decision-making processes and evaluating outcomes in future interactions, leveraging accumulated historical insights.

\subsection{CoT Prompting}

Addressing complex rationales necessitates that LLMs engage in extended chains of reasoning, a process distinct from the reasoning across disparate factual information typical of fact queries. However, Chain-of-Thoughts~\cite{wei2022chain}, Tree-of-Thoughts~\cite{yao2024tree} or Graph-of-Thoughts~\cite{yao2023beyond} methodologies proves effective for such scenarios. For issues that are well-studied and have high general applicability, manually designing CoT prompts emerges as a feasible solution. Ji et al. (2023)~\cite{ji2023towards} proposed a method of self-reflection that integrates knowledge acquisition with answer generation. By utilizing external tools and designing prompts, they constructed three types of self-reflection loops: the Factual Knowledge Acquiring Loop, the Knowledge-Consistent Answering Loop, and the Question-Entailment Answering Loop, thereby incorporating external rationales into the model's processing. Furthermore, Wu et al. (2024)~\cite{wu2024chain} conducted a manual analysis of error types in clinical records and developed three distinct CoT prompts to direct the GPT-4 model~\cite{achiam2023gpt} in focusing on intervention, diagnostic, and management errors. This targeted prompting facilitates the tasks of automatic error detection, span identification, and correction within clinical records.

While manual design of CoT prompts is highly effective, it requires substantial human and temporal resources. To mitigate these costs, Automate-CoT~\cite{shum2023automatic} proposed a technique for generating augmenting rational chains from a minimally labeled dataset. This approach employs a variance-reduced policy gradient strategy to evaluate the importance of each CoT chain, thus facilitating the selection of the most effective prompt combination.

Another form of utilizing Chain of Thoughts prompting involves constructing an agent workflow centered around LLMs. This typically requires the development of a more comprehensive system to address various real-world scenarios. According to Wang et al., such systems can be broadly divided into profiling, memory, planning, and action modules~\cite{wang2024survey}. Interpretable rationales can be integrated into multiple modules in various forms, allowing the agent to adapt and iterate based on environmental or human feedback. Recent advancements, such as those by LLM Reasoners~\cite{hao2024llm} and SocREval~\cite{he2024socreval}, have focused on automatically evaluating the quality of reasoning chains. These methodologies also assist in constructing robust \app{}.

Applications based on interpretable rationales span various domains. For instance, CoML~\cite{zhang-etal-2024-mlcopilot} integrates AutoML knowledge as prompts into an LLM, dynamically retrieves useful information from historical experimental records, and combines these elements to empower the LLM to develop machine learning solutions for novel tasks. MetaGPT~\cite{hong2023metagpt} has developed a multi-agent system for software development, where different stakeholders within a project are each represented as an agent. This setup enables multiple agents to collaborate according to a real-world work pipeline, effectively completing software development tasks. Similarly, sophisticated agent systems have been designed in fields such as customer service~\cite{abbasian2023conversational} and medical question answering~\cite{tang2023medagents}. In these domains, agents are tailored to handle specific types of inquiries, which can involve understanding complex user requests or providing accurate medical information. These systems not only enhance the interaction quality but also improve the efficiency and accuracy of responses, demonstrating the versatility and potential of LLMs when integrated into well-designed agent workflows.